\begin{exercises}
  \item 
    对于
    \begin{equation*}
      T=
      \begin{mat}[r]
        1  &3  \\
       -2  &-6
      \end{mat}
      \quad
      \hat{T}=
      \begin{mat}[r]
        0    &0  \\
       -11/2 &-5
      \end{mat}
      \quad
      P=
      \begin{mat}[r]
        4  &2  \\
       -3  &2
      \end{mat}
    \end{equation*}
    验证 $\hat{T}=PTP^{-1}$。
    \begin{answer}
      一种做法是从左到右计算。
      \begin{multline*}
        PTP^{-1}=
        \begin{mat}[r]
          4  &2  \\
         -3  &2
        \end{mat}
        \begin{mat}[r]
          1  &3  \\
         -2  &-6
        \end{mat}
        \begin{mat}[r]
          2/14  &-2/14  \\
          3/14  &4/14
        \end{mat}                                 \\
        =
        \begin{mat}[r]
          0  &0  \\
         -7  &-21
        \end{mat}  
        \begin{mat}[r]
          2/14  &-2/14  \\
          3/14  &4/14
        \end{mat}
        =
        \begin{mat}[r]
          0    &0  \\
         -11/2 &-5
        \end{mat}
      \end{multline*}
    \end{answer}
  \item
    \nearbyexample{ex:OnlyZeroSimToZero} 表明，与零矩阵相似的矩阵只有它本身，
    而且
    与单位矩阵相似的矩阵
    也只有它本身。
    \begin{exparts}
      \partsitem 证明：唯一元素为 $2$ 的 $\nbyn{1}$ 矩阵
         也只与它自身相似。
      \partsitem 形如 $cI$（$c$ 为某个标量）的矩阵
         是否也只与它自身相似？
     \partsitem 对角矩阵是否只与它自身相似？
     % \partsitem Is a square block partial-identity matrix similar only to 
     % itself?
    \end{exparts}
    \begin{answer}
     \begin{exparts}
      \partsitem 由于矩阵 $(2)$ 是 $\nbyn{1}$ 的，矩阵
         $P$ 与 $P^{-1}$ 也是 $\nbyn{1}$ 的，因此当
         $P=(p)$ 时，其逆为 $P^{-1}=(1/p)$。  
         于是 $P(2)P^{-1}=(p)(2)(1/p)=(2)$。
       \partsitem 是的：~回忆一下，我们可以把标量倍数提到矩阵外面 
        \( P(cI)P^{-1}=cPIP^{-1}=cI \)。
        顺便指出，零矩阵与单位矩阵分别是特殊情形
        $c=0$ 与 $c=1$。
      \partsitem 不是，如下面这个例子所示。
        \begin{equation*}
           \begin{mat}[r]
              1  &-2  \\
             -1  &1
            \end{mat}
           \begin{mat}[r]
             -1  &0   \\
              0  &-3
           \end{mat}
           \begin{mat}[r]
              -1  &-2   \\
              -1  &-1
           \end{mat}
           =
           \begin{mat}[r]
              -5  &-4   \\
              2   &1
           \end{mat}
        \end{equation*} 
    \end{exparts}  
   \end{answer}
  \recommended \item 考虑~$\C^3$ 的这个变换
    \begin{equation*}
      t(\colvec{x \\ y \\ z})=\colvec{x-z \\ z \\ 2y}
    \end{equation*}
    以及这些基。
    \begin{equation*}
      B=\sequence{\colvec{1 \\ 2 \\ 3}, 
                  \colvec{0 \\ 1 \\ 0}, 
                  \colvec{0 \\ 0 \\ 1}}
      \qquad
      D=\sequence{\colvec{1 \\ 0 \\ 0},
                  \colvec{1 \\ 1 \\ 0},
                  \colvec{1 \\ 0 \\ 1}}
    \end{equation*}
    我们将计算箭头图的各个部分，
    以使用两个相似矩阵来表示该变换。
    \begin{exparts}
      \partsitem 画出针对本情形具体化的箭头图。
      \partsitem 计算 $T=\rep{t}{B,B}$。
      \partsitem 计算 $\hat{T}=\rep{t}{D,D}$。
      \partsitem 计算箭头正方形另外两条边的
         矩阵。
    \end{exparts}
    \begin{answer}
      \begin{exparts}
        \partsitem
          \begin{equation*}
            \begin{CD}
              \C^3_{\wrt{B}}                   @>t>T>        \C^3_{\wrt{B}}       \\
              @V{\scriptstyle\identity} VV              @V{\scriptstyle\identity} VV \\
              \C^3_{\wrt{D}}                   @>t>\hat{T}>        \C^3_{\wrt{D}}
            \end{CD}
          \end{equation*}
        \partsitem
          对于起始基~$B$ 中的每个元素，求 
          变换作用于其后的结果
          \begin{equation*}
             \colvec{1 \\ 2 \\ 3}\mapsunder{t}\colvec{-2 \\ 3 \\ 4}
             \qquad 
             \colvec{0 \\ 1 \\ 0}\mapsunder{t}\colvec{0 \\ 0 \\ 2}
             \qquad 
             \colvec{0 \\ 0 \\ 1}\mapsunder{t}\colvec{-1 \\ 1 \\ 0}
          \end{equation*}
          并将这些输出关于终止基~$B$ 表示
          \begin{equation*}
            \rep{\colvec{-2 \\ 3 \\ 4}}{B}=\colvec{-2 \\ 7 \\ 10}
            \qquad
            \rep{\colvec{0 \\ 0 \\ 2}}{B}=\colvec{0 \\ 0 \\ 2}
            \qquad
            \rep{\colvec{-1 \\ 1 \\ 0}}{B}=\colvec{-1 \\ 3 \\ 3}
          \end{equation*}
          由此得到矩阵。
          \begin{equation*}
            T=\rep{t}{B,B}=
            \begin{mat}
              -2 &0 &-1 \\
               7 &0 &3  \\
              10 &2 &3 
            \end{mat}
          \end{equation*}
        \partsitem
          求变换作用于~$D$ 中元素的结果
          \begin{equation*}
             \colvec{1 \\ 0 \\ 0}\mapsunder{t}\colvec{1 \\ 0 \\ 0}
             \qquad 
             \colvec{1 \\ 1 \\ 0}\mapsunder{t}\colvec{1 \\ 0 \\ 2}
             \qquad 
             \colvec{1 \\ 0 \\ 1}\mapsunder{t}\colvec{0 \\ 1 \\ 0}
          \end{equation*}
          并将它们关于终止基~$D$ 表示
          \begin{equation*}
            \rep{\colvec{1 \\ 0 \\ 0}}{D}=\colvec{1 \\ 0 \\ 0}
            \qquad
            \rep{\colvec{1 \\ 0 \\ 2}}{D}=\colvec{-1 \\ 0 \\ 2}
            \qquad
            \rep{\colvec{0 \\ 1 \\ 0}}{D}=\colvec{-1 \\ 1 \\ 0}
          \end{equation*}
          由此得到矩阵。
          \begin{equation*}
            \hat{T}=\rep{t}{D,D}=
            \begin{mat}
               1 &-1 &-1 \\
               0 &0  &1  \\
               0 &2  &0 
            \end{mat}
          \end{equation*}
        \partsitem
          要沿右侧向下，我们需要 
          $\rep{\identity}{B,D}$
          因此我们先计算恒等映射对 
          ~$B$ 的每个元素的作用——它没有任何作用，
          然后将结果关于
          ~$D$ 表示。 
          \begin{equation*}
            \rep{\colvec{1 \\ 2 \\ 3}}{D}=\colvec{-4 \\ 2 \\ 3}
            \qquad
            \rep{\colvec{0 \\ 1 \\ 0}}{D}=\colvec{-1 \\ 1 \\ 0}
            \qquad
            \rep{\colvec{0 \\ 0 \\ 1}}{D}=\colvec{-1 \\ 0 \\ 1}
          \end{equation*}
          所以这就是~$P$。
          \begin{equation*}
            P=
            \begin{mat}
              -4 &-1 &-1 \\
               2 &1  &0  \\
               3 &0  &1
            \end{mat}
          \end{equation*}
          对于另一个矩阵~$\rep{\identity}{D,B}$，我们既可以 
          像刚才对~$P$ 那样直接求出它，也可以进行 
          常规的矩阵求逆计算。
          \begin{equation*}
            P^{-1}=
            \begin{mat}
               1 &1 &1 \\
               -2 &-1  &-2  \\
               -3 &-3  &-2
            \end{mat}
          \end{equation*}
      \end{exparts}
    \end{answer}
  \item 
    考虑变换 $\map{t}{\polyspace_2}{\polyspace_2}$
    它由
    $x^2\mapsto x+1$、$x\mapsto x^2-1$ 以及 $1\mapsto 3$ 描述。
    \begin{exparts}
      \partsitem 求 $T=\rep{t}{B,B}$，其中 $B=\sequence{x^2,x,1}$。 
      \partsitem 求 $\hat{T}=\rep{t}{D,D}$，其中 $D=\sequence{1,1+x,1+x+x^2}$。
      \partsitem 求满足 $\hat{T}=PTP^{-1}$ 的矩阵 $P$。 
    \end{exparts}
    \begin{answer}
      \begin{exparts}
        \partsitem 由于我们用 $B$ 的成员来描述 $t$，
          求矩阵表示很容易：
          \begin{equation*}
            \rep{t(x^2)}{B}=\colvec[r]{0 \\ 1 \\ 1}_B
            \quad
            \rep{t(x)}{B}=\colvec[r]{1 \\ 0 \\ -1}_B
            \quad
            \rep{t(1)}{B}=\colvec[r]{0 \\ 0 \\ 3}_B
          \end{equation*}
          由此得到下式。
          \begin{equation*}
            \rep{t}{B,B}
            \begin{mat}[r]
              0  &1  &0  \\
              1  &0  &0  \\
              1  &-1 &3  
            \end{mat}
          \end{equation*}
        \partsitem 我们将求 $t(1)$、$t(1+x)$ 和 $t(1+x+x^2$，
          以求出每一个关于 $D$ 的表示。
          已知 $t(1)=3$，另外两个容易看出：
          $t(1+x)=x^2+2$ 以及 $t(1+x+x^2)=x^2+x+3$。
          通过观察，我们得到每个向量的表示
          \begin{equation*}
            \rep{t(1)}{D}=\colvec[r]{3 \\ 0 \\ 0}_D
            \quad
            \rep{t(1+x)}{D}=\colvec[r]{2  \\ -1 \\  1}_D
            \quad
            \rep{t(1+x+x^2)}{D}=\colvec[r]{2 \\ 0 \\ 1}_D
          \end{equation*}
          从而得到该映射的表示。
          \begin{equation*}
            \rep{t}{D,D}
            =
            \begin{mat}[r]
              3  &2  &2  \\
              0  &-1 &0  \\
              0  &1  &1
            \end{mat}
          \end{equation*}
         \partsitem 下面的图表
           \begin{equation*}
             \begin{CD}
               V_{\wrt{B}}                  @>t>T>  V_{\wrt{B}}       \\
               @V\scriptstyle\identity VPV      @V\scriptstyle\identity VPV \\
               V_{\wrt{D}}                  @>t>\hat{T}>  V_{\wrt{D}}
             \end{CD}
           \end{equation*}
           表明它们是 $P=\rep{\identity}{B,D}$ 与 
           $P^{-1}=\rep{\identity}{D,B}$。
           \begin{equation*}
             P=
             \begin{mat}[r]
               0  &-1  &1  \\ 
               -1  &1  &0  \\
               1  &0  &0 
             \end{mat}
             \qquad
             P^{-1}=
             \begin{mat}[r]
               0  &0  &1  \\ 
               0  &1  &1  \\
               1  &1  &1 
             \end{mat}
           \end{equation*}
      \end{exparts}
    \end{answer}
  \recommended \item 
    设 $T$ 是 $\map{t}{\C^2}{\C^2}$ 关于 $B,B$ 的表示。
    \begin{equation*}
      T=
      \begin{mat}
        1 &-1 \\ 
        2 &1
      \end{mat}
      \qquad
      B=\sequence{\colvec{1 \\ 0},\colvec{1 \\ 1}},\hspace{0.7em}
      D=\sequence{\colvec{2 \\ 0},\colvec{0 \\ -2}}
    \end{equation*}
    我们将转换为关于 $D,D$ 表示~$t$ 的矩阵。
    \begin{exparts}
      \partsitem 画出箭头图。
      \partsitem 给出表示该图左边与右边
          两条边的矩阵，方向按我们为完成转换
          而遍历该图的方向。
      \partsitem 求 $\rep{t}{D,D}$。
    \end{exparts}
    \begin{answer}
      \begin{exparts}
        \partsitem
          \begin{equation*}
            \begin{CD}
              \C^2_{\wrt{B}}                   @>t>T>        \C^2_{\wrt{B}}       \\
              @V{\scriptstyle\identity} VV              @V{\scriptstyle\identity} VV \\
              \C^2_{\wrt{D}}                   @>t>\hat{T}>        \C^2_{\wrt{D}}
            \end{CD}
          \end{equation*}  
        \partsitem
          对于右边，我们求恒等映射的作用，
          它没有任何作用，
          \begin{equation*}
            \colvec{1 \\ 0}\mapsunder{\identity}\colvec{1 \\ 0}
            \qquad
            \colvec{1 \\ 1}\mapsunder{\identity}\colvec{1 \\ 1}
          \end{equation*}
          并将它们关于~$D$ 表示
          \begin{equation*}
            \rep{\colvec{1 \\ 0}}{D}=\colvec{1/2 \\ 0}
            \qquad
            \rep{\colvec{1 \\ 1}}{D}=\colvec{1/2 \\ -1/2}
          \end{equation*}
          于是我们得到下式。
          \begin{equation*}
            P=\rep{\identity}{B,D}=
            \begin{mat}
              1/2 &1/2 \\
              0   &-1/2
            \end{mat}
          \end{equation*}

          对于左边的矩阵，我们既可以像 
          上一段那样直接计算，
          也可以取其逆。 
          \begin{equation*}
            P^{-1}=\rep{\identity}{D,B}=
            \frac{1}{(-1/4)}\cdot
            \begin{mat}
              -1/2 &-1/2 \\
              0   &1/2
            \end{mat}
            =
            \begin{mat}
              2 &2 \\
              0 &-2
            \end{mat}
          \end{equation*}
        \partsitem
          与前一问一样，我们既可以 
          由定义直接计算，
          也可以利用矩阵运算来计算。
          \begin{equation*}
            PTP^{-1}=
            \begin{mat}
              2 &2 \\
              0 &-2
            \end{mat}
            \begin{mat}
              1 &-1 \\ 
              2 &1
            \end{mat}
            \begin{mat}
              2 &2 \\
              0 &-2
            \end{mat}
            =
            \begin{mat}
              3 &3  \\
             -2 &-1
            \end{mat}
          \end{equation*}
      \end{exparts}
    \end{answer}
  \recommended \item
     通过让
     \( \map{t}{\C^2}{\C^2} \) 按下面的方式作用，给出一个非平凡的相似关系，
     \begin{equation*}
        \colvec[r]{1 \\ 2}\mapsto\colvec[r]{3 \\ 0}
        \qquad
        \colvec[r]{-1 \\ 1}\mapsto\colvec[r]{-1 \\ 2}
     \end{equation*}
     选取两个基~$B,D$，
     并关于它们表示 \( t \)，
     即 \( \hat{T}=\rep{t}{B,B} \) 与 \( T=\rep{t}{D,D} \)。
     然后计算 
     用于把基从 \( B \) 换到 \( D \) 以及 
     换回来的 \( P \) 与 \( P^{-1} \)。
     \begin{answer}
        基的一种可行选取是 
        \begin{equation*}
          B=\sequence{\colvec[r]{1 \\ 2},\colvec[r]{-1 \\ 1}}
          \qquad
          D=\stdbasis_2=\sequence{\colvec[r]{1 \\ 0},\colvec[r]{0 \\ 1}}
        \end{equation*}
        （这个 $B$ 来自对该映射的描述）。
        为求矩阵 $\hat{T}=\rep{t}{B,B}$，解关系式 
        \begin{equation*}
           c_1\colvec[r]{1 \\ 2}+c_2\colvec[r]{-1 \\ 1}=\colvec[r]{3 \\ 0}
           \qquad
          \hat{c}_1\colvec[r]{1 \\ 2}+\hat{c}_2\colvec[r]{-1 \\ 1}=\colvec[r]{-1 \\ 2}
        \end{equation*}
        得 \( c_1=1 \)、\( c_2=-2 \)、\( \hat{c}_1=1/3 \) 及
        \( \hat{c}_2=4/3 \)。
        \begin{equation*}
           \rep{t}{B,B}=
           \begin{mat}[r]
              1  &1/3 \\
             -2  &4/3
           \end{mat}
        \end{equation*}

        求 \( \rep{t}{D,D} \) 需要稍多一些计算。
        我们先求 \( t(\vec{e}_1) \)。
        关系式
        \begin{equation*}
          c_1\colvec[r]{1 \\ 2}+c_2\colvec[r]{-1 \\ 1}=\colvec[r]{1 \\ 0}
        \end{equation*}
        给出 \( c_1=1/3 \) 与 \( c_2=-2/3 \)，于是 
        \begin{equation*}
          \rep{\vec{e}_1}{B}=\colvec[r]{1/3 \\ -2/3}_B
        \end{equation*}
        从而
        \begin{equation*}
           \rep{t(\vec{e}_1)}{B}=
                  \begin{mat}[r]
                     1  &1/3  \\
                    -2  &4/3
                  \end{mat}_{B,B}
                  \colvec[r]{1/3 \\ -2/3}_B
                  =
                  \colvec[r]{1/9 \\ -14/9}_B
        \end{equation*}
        因此 $t$ 对第一个基向量 $\vec{e}_1$ 的作用如下。
        \begin{equation*}
          t(\vec{e}_1)
          =(1/9)\cdot\colvec[r]{1 \\ 2}-(14/9)\cdot\colvec[r]{-1 \\ 1}
          =\colvec[r]{5/3 \\ -4/3}      
        \end{equation*}
        \( t(\vec{e}_2) \) 的计算类似。
        关系式
        \begin{equation*}
          c_1\colvec[r]{1 \\ 2}+c_2\colvec[r]{-1 \\ 1}=\colvec[r]{0 \\ 1}
        \end{equation*}
        给出 \( c_1=1/3 \) 与 \( c_2=1/3 \)，所以
        \begin{equation*}
          \rep{\vec{e}_1}{B}=\colvec[r]{1/3 \\ 1/3}_B      
        \end{equation*}
        从而
        \begin{equation*}
          \rep{t(\vec{e}_1)}{B}=
                  \begin{mat}[r]
                     1  &1/3  \\
                    -2  &4/3
                  \end{mat}_{B,B}
                  \colvec[r]{1/3 \\ 1/3}_B
                  =
                  \colvec[r]{4/9 \\ -2/9}_B
        \end{equation*}
        因此 $t$ 对第二个基向量 $\vec{e}_2$ 的作用如下。
        \begin{equation*}
          t(\vec{e}_2)
          =(4/9)\cdot\colvec[r]{1 \\ 2}-(2/9)\cdot\colvec[r]{-1 \\ 1}
          =\colvec[r]{2/3 \\ 2/3}
        \end{equation*}
        因此
        \begin{equation*}
           \rep{t}{D,D}=
           \begin{mat}[r]
              5/3  &2/3  \\
             -4/3  &2/3
           \end{mat}
        \end{equation*}
        于是得到这个矩阵。
        \begin{equation*}
          P=\rep{\identity}{B,D}
          =\begin{mat}[r]
            1  &-1  \\
            2  &1 
          \end{mat}
        \end{equation*}
       而这一个用于换基。 
       \begin{equation*}
          P^{-1}=\bigl(\rep{\identity}{B,D}\bigr)^{-1}
          =\begin{mat}[r]
             1  &-1 \\
             2  &1
          \end{mat}^{-1}
          =
          \begin{mat}[r]
            1/3  &1/3  \\
            -2/3 &1/3
          \end{mat}
       \end{equation*}
       这些计算的检验是例行的。
       \begin{equation*}
          \begin{mat}[r]
             1  &-1 \\
             2  &1
          \end{mat}
          \begin{mat}[r]
             1  &1/3 \\
            -2  &4/3
          \end{mat}
          \begin{mat}[r]
            1/3 &1/3 \\
           -2/3 &1/3
          \end{mat}
          =
          \begin{mat}[r]
            5/3 &2/3 \\
           -4/3 &2/3
          \end{mat}
       \end{equation*}  
    \end{answer}
  \recommended \item 
    证明这两个矩阵不相似。
    \begin{equation*}
       \begin{mat}[r]
          1  &0  &4  \\
          1  &1  &3  \\
          2  &1  &7
       \end{mat}
       \qquad
       \begin{mat}[r]
          1  &0  &1  \\
          0  &1  &1  \\
          3  &1  &2
       \end{mat}
    \end{equation*}
    \begin{answer}
      高斯消元法表明，
      第一个矩阵表示的映射秩为二，而第二个
      矩阵表示的映射秩为三。
    \end{answer}
  \item 
    用映射的语言解释 \nearbyexample{ex:OnlyZeroSimToZero}。
    \begin{answer}
      零映射的表示只有零矩阵，
      无论基对是什么都有 $\rep{z}{B,D}=Z$，
      因此特别地，对任何单独一个基 $B$ 我们有 $\rep{z}{B,B}=Z$。  
      恒等映射的情形稍有不同：~恒等映射关于任何 $B,B$ 的
      唯一表示
      是单位矩阵 $\rep{\identity}{B,B}=I$。
      （\textit{注：}~当然，我们已经见过 $B\neq D$ 且 
      $\rep{\identity}{B,D}\neq I$ 的例子 \Dash 事实上，我们已经看到任何 
      非奇异矩阵都是恒等映射关于某个 
      $B,D$ 的表示。）
    \end{answer}
  \recommended \item 
    \cite{Halmos}
    是否存在两个矩阵 \( A \) 与 \( B \) 相似，
    而 \( A^2 \) 与 \( B^2 \) 却不相似？
    \begin{answer}
      不存在。
      若 \( A=PBP^{-1} \)，则 \( A^2=(PBP^{-1})(PBP^{-1})=PB^2P^{-1} \)。
    \end{answer}
  \recommended \item
    证明：若两个矩阵相似且其中一个是可逆的，
    则另一个也可逆。
    \begin{answer}
       矩阵相似是矩阵等价的特殊情形 
       （若矩阵相似则它们矩阵等价），
       而矩阵等价保持非奇异性。
    \end{answer}
  \item \label{exer:SimIsEquivRel}
    证明相似是一个等价关系。
    （前面给出的定义已经体现了这一点，因此 
    这里改从如下定义出发：$\hat{T}$ 相似于
    $T$ 当且仅当 $\hat{T}=PTP^{-1}$。）
    \begin{answer}
       矩阵与它自身相似；取 \( P \) 为单位
       矩阵：~$P=IPI^{-1}=IPI$。

       若 \( \hat{T} \) 相似于 \( T \)，则 \( \hat{T}=PTP^{-1} \)
       于是 \( P^{-1}\hat{T}P=T \)。
       将其改写为 \( T=(P^{-1})\hat{T}(P^{-1})^{-1} \)，即可得出
       $T$ 相似于 \( \hat{T} \)。

        对于传递性，
        若 \( T \) 相似于 \( S \) 且 \( S \) 相似于 \( U \)，
        则 \( T=PSP^{-1} \) 且 \( S=QUQ^{-1} \)。
        于是 \( T=PQUQ^{-1}P^{-1}=(PQ)U(PQ)^{-1} \)，这表明 \( T \)
        相似于 \( U \)。  
      \end{answer}
  \item 
      考虑一个 
      关于某个 $B,B$ 表示 
      \( \Re^2 \) 中关于 \( x \) 轴的反射的矩阵。
      再考虑  
      一个关于某个 $D,D$ 表示
      关于 \( y \) 轴的反射的矩阵。
      它们必定相似吗？
      \begin{answer}
         设 $f_x$ 与 $f_y$ 为这两个反射映射（有时称为「翻转」）。
         对任何基
         \( B \) 与 \( D \)，矩阵 \( \rep{f_x}{B,B}  \) 与
         \( \rep{f_y}{D,D} \) 相似。
         首先注意
         \begin{equation*}
            S=\rep{f_x}{\stdbasis_2,\stdbasis_2}=
            \begin{mat}[r]
               1  &0  \\
               0  &-1
            \end{mat}
            \qquad
            T=\rep{f_y}{\stdbasis_2,\stdbasis_2}=
            \begin{mat}[r]
              -1  &0  \\
               0  &1
            \end{mat}
         \end{equation*}
         相似，因为第二个矩阵是 $f_x$
         关于基 \( A=\sequence{\vec{e}_2,\vec{e}_1} \) 的表示：
         \begin{equation*}
            \begin{mat}[r]
               1  &0  \\
               0  &-1
            \end{mat}
            =    
            P
            \begin{mat}[r]
              -1  &0  \\
               0  &1
            \end{mat}
            P^{-1}
         \end{equation*}
         其中 $P=\rep{\identity}{A,\stdbasis_2}$。
         \begin{equation*}
           \begin{CD}
             \C^2_{\wrt{A}}                   
                @>f_x>T>        
                V\C^2_{\wrt{A}}       \\
             @V\scriptstyle\identity VPV                
                @V\scriptstyle\identity VPV \\
             \C^2_{\wrt{\stdbasis_2}}         
                @>f_x>S>        
                \C^2_{\wrt{\stdbasis_2}}
           \end{CD}
         \end{equation*}
         现在结论由
         \nearbyexercise{exer:SimIsEquivRel} 的传递性部分得出。
         
        我们也可以不依赖那道练习来完成。
        记
        $\rep{f_x}{B,B}=QTQ^{-1}=Q\rep{f_x}{\stdbasis_2,\stdbasis_2}Q^{-1}$
        以及
        $\rep{f_y}{D,D}=RSR^{-1}=R\rep{f_y}{\stdbasis_2,\stdbasis_2}R^{-1}$。
        由第一段中的等式， 
        这两个中的第一个是
        $\rep{f_x}{B,B}=QP\rep{f_y}{\stdbasis_2,\stdbasis_2}P^{-1}Q^{-1}$。
        将这两个中的第二个改写为
        $R^{-1}\cdot\rep{f_y}{D,D}\cdot R=\rep{f_y}{\stdbasis_2,\stdbasis_2}$
        并代入，就得到所需的关系
        \begin{multline*}
          \rep{f_x}{B,B}
          =QP\rep{f_y}{\stdbasis_2,\stdbasis_2}P^{-1}Q^{-1}  \\
          =QPR^{-1}\cdot\rep{f_y}{D,D}\cdot RP^{-1}Q^{-1}
          =(QPR^{-1})\cdot\rep{f_y}{D,D}\cdot (QPR^{-1})^{-1}
        \end{multline*}
        因此矩阵 \( \rep{f_x}{B,B}  \) 与 \( \rep{f_y}{D,D} \) 
        相似。 
      \end{answer}
  \item 
     证明矩阵相似保持秩与行列式。
     逆命题成立吗？
     \begin{answer}
        我们必须证明：若两个矩阵相似，则它们有相同的
        秩和相同的行列式。

        首先，我们已经证明 
        映射的秩等于表示该映射的任何矩阵的秩，
        所以秩对所有表示都相同。
        （换句话说，秩是映射的性质， 
        因为它是值域空间的维数。）

        至于行列式，
        \( \deter{PSP^{-1}}=\deter{P}\cdot\deter{S}\cdot\deter{P^{-1}}
           =\deter{P}\cdot\deter{S}\cdot\deter{P}^{-1}=\deter{S} \)。

       该命题的逆命题不成立。 
       一方面，
       存在行列式相同却不相似的矩阵。
       举个例子，考虑一个行列式为 
       零的非零矩阵。
       它不与零矩阵相似，因为零矩阵只与
       它自身相似，但二者行列式相同。 
       同样的思路给出秩的反例。 
     \end{answer}
  \item 
     是否存在只包含一个矩阵相似
     类的矩阵等价类？
     是否存在包含无穷多个相似类的？
     \begin{answer}
       包含所有秩为 
       零的 \( \nbyn{n} \) 矩阵的矩阵等价类只含一个矩阵，即零矩阵。
       因此它作为子集只含有一个相似类。

       相比之下，秩为一的 
       \( \nbyn{1} \) 矩阵的矩阵等价类由那些 
       \( k\neq 0 \) 的 $\nbyn{1}$ 矩阵 \( (k) \) 组成。
       对任何基 \( B \)，
       用标量 \( k \) 相乘的映射的表示
       为 \( \rep{t_k}{B,B}=(k) \)，
       所以每个这样的矩阵独自构成其相似类。
       所以这是矩阵等价类分裂成
       无穷多个相似类的一个例子。  
      \end{answer}
  \item 
     两个不同的对角矩阵能属于同一个相似类吗？
     \begin{answer}
       能，下面两个矩阵相似
       \begin{equation*}
          \begin{mat}[r]
            1  &0  \\
            0  &3
          \end{mat}
          \qquad
          \begin{mat}[r]
            3  &0  \\
            0  &1
          \end{mat}
       \end{equation*}
       因为，当第一个矩阵是关于 
       $B=\sequence{\vec{\beta}_1,\vec{\beta}_2}$ 的 $\rep{t}{B,B}$ 时， 
       第二个矩阵就是关于 
       $D=\sequence{\vec{\beta}_2,\vec{\beta}_1}$ 的 $\rep{t}{D,D}$。
      \end{answer}
  \recommended \item
     证明：若两个矩阵相似，则当 \( k>0 \) 时，
     它们的 \( k \) 次幂也相似。
     若 \( k\leq 0 \) 会怎样？
     \begin{answer}
       \( k \) 次幂相似，因为当每个矩阵表示
       映射 $t$ 时，\( k \) 次幂表示
       \( t^k \)，即 $k$ 个 $t$ 的复合。
       （例如，若 $T=rep{t}{B,B}$，则 $T^2=\rep{\composed{t}{t}}{B,B}$。）

       用更计算化的说法重述：若 \( T=PSP^{-1} \)，则
       \( T^2=(PSP^{-1})(PSP^{-1})=PS^2P^{-1} \)。
       归纳法将此推广到所有幂。

       对于 $k\leq 0$ 的情形， 
       设 \( S \) 可逆且 \( T=PSP^{-1} \)。
       注意 \( T \) 可逆：
       \( T^{-1}=(PSP^{-1})^{-1}=PS^{-1}P^{-1} \)，
       而同一个等式表明 
       \( T^{-1} \) 相似于 \( S^{-1} \)。
       其余负次幂则由第一段给出。  
     \end{answer}
  \recommended \item
     设 \( p(x) \) 为多项式 \( c_nx^n+\cdots+c_1x+c_0 \)。
     证明：若 \( T \) 相似于 \( S \)，则
     \( p(T)=c_nT^n+\cdots+c_1T+c_0I \) 相似于
     \( p(S)=c_nS^n+\cdots+c_1S+c_0I \)。
     \begin{answer}
        从概念上说，二者都表示某个
        变换 \( t \) 的 \( p(t) \)。
        从计算上说，我们有下式。
        \begin{align*}
           p(T)
           &=c_n(PSP^{-1})^n+\dots+c_1(PSP^{-1})+c_0I   \\
           &=c_nPS^nP^{-1}+\dots+c_1PSP^{-1}+c_0I   \\
           &=Pc_nS^nP^{-1}+\dots+Pc_1SP^{-1}+Pc_0P^{-1}   \\
           &=P(c_nS^n+\dots+c_1S+c_0)P^{-1}
        \end{align*}  
      \end{answer}
  \item 
     列出 \( \nbyn{1} \) 矩阵的所有矩阵等价类。
     同时列出相似类，并描述每个矩阵等价类
     内部包含哪些相似类。
     \begin{answer}
       有两个等价类：(i)~秩为零的矩阵的类， 
       其中只有一个矩阵：
       $\mathscr{C}_1=\set{(0)}$，
       以及 (2)~秩为一的矩阵的类，
       其中有无穷多个： 
       \( \mathscr{C}_2=\set{(k)\suchthat k\neq0} \)。

      每个 \( \nbyn{1} \) 矩阵都独自构成其相似类。
      这是因为一维空间上的任何变换
      都是乘以一个标量的映射 $\map{t_k}{V}{V}$，由
      $\vec{v}\mapsto k\cdot\vec{v}$ 给出。 
      于是，对任何基 \( B=\sequence{\vec{\beta}} \)，
      表示变换 \( t_k \) 的矩阵 
      关于 \( B,B \) 为
      \( (\rep{t_k(\vec{\beta})}{B})=(k) \)。 
      
      所以，矩阵等价类
      $\mathscr{C}_1$ 中（显然）只含有由矩阵 $(0)$ 构成的 
      唯一相似类。
      而矩阵等价类 $\mathscr{C}_2$ 中含有
      无穷多个各含一个成员的相似类，由
      $k\neq0$ 的 $(k)$ 构成。  
    \end{answer}
  \item 
     相似保持和吗？
     \begin{answer}
       不保持。
       下面是一个例子，其中有两对矩阵，每对中的两个矩阵相似：
       \begin{equation*}
          \begin{mat}[r]
             1  &-1  \\
             1  &2
          \end{mat}
          \begin{mat}[r]
             1  &0   \\
             0  &3
          \end{mat}
          \begin{mat}[r]
             2/3  &1/3   \\
            -1/3  &1/3
          \end{mat}
          =
          \begin{mat}[r]
             5/3  &-2/3  \\
            -4/3  &7/3
          \end{mat}
       \end{equation*}
       以及
       \begin{equation*}
          \begin{mat}[r]
             1  &-2  \\
            -1  &1
          \end{mat}
          \begin{mat}[r]
            -1  &0   \\
             0  &-3
          \end{mat}
          \begin{mat}[r]
             -1  &-2   \\
             -1  &-1
          \end{mat}
          =
          \begin{mat}[r]
             -5  &-4   \\
              2  &1
          \end{mat}
       \end{equation*}
       （这个例子并非完全任意，因为
       两个等号左边的中间矩阵相加为零矩阵）。
       注意这些相似矩阵的和并不相似
       \begin{equation*}
          \begin{mat}[r]
             1  &0   \\
             0  &3
          \end{mat}
          +
          \begin{mat}[r]
            -1  &0   \\
             0  &-3
          \end{mat}
          =
          \begin{mat}[r]
            0  &0  \\
            0  &0
          \end{mat}
          \qquad
          \begin{mat}[r]
             5/3  &-2/3   \\
             -4/3 &7/3
          \end{mat}
          +
          \begin{mat}[r]
             -5  &-4   \\
              2  &1
          \end{mat}
          \neq
          \begin{mat}[r]
            0  &0  \\
            0  &0
          \end{mat}
       \end{equation*}
       因为零矩阵只与它自身相似。
     \end{answer}
  \item 
     证明：若 \( T-\lambda I \) 与 \( N \) 是相似矩阵，则
     \( T \) 与 \( N+\lambda I \) 也相似。
     \begin{answer}
    若 \( N=P(T-\lambda I)P^{-1} \)，则
    \( N=PTP^{-1}-P(\lambda I)P^{-1} \)。
    对角矩阵 \( \lambda I \) 与任何矩阵可交换，所以
    \( P(\lambda I)P^{-1}=PP^{-1}(\lambda I)=\lambda I \)。
    于是 \( N=PTP^{-1}-\lambda I \)，
    从而 \( N+\lambda I=PTP^{-1} \)。
    （所以二者不仅相似，事实上它们是经由 
    同一个 \( P \) 相似的。）
    \end{answer}
\end{exercises}
